\documentclass[11pt,a4paper]{article}

\usepackage{amsmath,amssymb}
\usepackage[margin=2.5cm]{geometry}

\begin{document}

\section*{Exercise 1: Parseval's Theorem (2 points)}

For the unitary DFT,
\[
F(\omega_k)=\frac{1}{\sqrt{N}}\sum_{x=0}^{N-1}f(x)e^{-i\omega_kx},
\qquad \omega_k=\frac{2\pi k}{N},
\]
prove that
\[
\sum_{x=0}^{N-1}|f(x)|^2=\sum_{k=0}^{N-1}|F(\omega_k)|^2.
\]

\noindent\textit{Hint:} Recall that $|z|^2=z\overline{z}$ and that
$\sum_{k=0}^{N-1}r^k=(1-r^N)/(1-r)$ for $r\neq 1$.

\subsection*{Your proof}

\textbf{Step 1: Write the left side using the inverse DFT.}
Since $|f(x)|^2 = f(x)\overline{f(x)}$, substitute the inverse DFT
$f(x) = \frac{1}{\sqrt{N}}\sum_{k=0}^{N-1} F(\omega_k) e^{i\omega_k x}$ into both
$f(x)$ and its conjugate (using independent summation indices $k$ and $k'$ so the two
sums can be kept apart):
\[
\sum_{x=0}^{N-1} |f(x)|^2
= \sum_{x=0}^{N-1} f(x)\overline{f(x)}
= \sum_{x=0}^{N-1} \left(\frac{1}{\sqrt{N}}\sum_{k=0}^{N-1} F(\omega_k) e^{i\omega_k x}\right)
\overline{\left(\frac{1}{\sqrt{N}}\sum_{k'=0}^{N-1} F(\omega_{k'}) e^{i\omega_{k'} x}\right)}.
\]

\textbf{Step 2: Expand the conjugate and rearrange the sums.}
Conjugating the second factor gives $\overline{F(\omega_{k'})}\,e^{-i\omega_{k'}x}$. Pulling
the $\frac{1}{\sqrt{N}}\cdot\frac{1}{\sqrt{N}} = \frac{1}{N}$ factor out and swapping the
(finite) sums over $x$, $k$, and $k'$ is valid since all sums are finite:
\[
\sum_{x=0}^{N-1} |f(x)|^2
= \frac{1}{N}\sum_{k=0}^{N-1}\sum_{k'=0}^{N-1} F(\omega_k)\,\overline{F(\omega_{k'})}
\underbrace{\sum_{x=0}^{N-1} e^{i(\omega_k-\omega_{k'})x}}_{=:S(k,k')}.
\]

\textbf{Step 3: Evaluate the orthogonality sum $S(k,k')$.}
Let $r = e^{i(\omega_k-\omega_{k'})} = e^{i2\pi(k-k')/N}$, so $S(k,k') = \sum_{x=0}^{N-1} r^x$.

\emph{Case $k=k'$:} then $r=1$, and the sum is just $N$ copies of $1$, so $S(k,k)=N$.

\emph{Case $k\neq k'$:} then $k-k'$ is a nonzero integer with $|k-k'|<N$, so
$\omega_k-\omega_{k'}=2\pi(k-k')/N$ is not a multiple of $2\pi$, hence $r\neq 1$ and the
geometric series formula from the hint applies:
\[
S(k,k') = \frac{1-r^N}{1-r}.
\]
But $r^N = e^{i2\pi(k-k')} = 1$ because $k-k'$ is an integer, so the numerator vanishes and
$S(k,k')=0$.

Combining both cases gives the orthogonality identity
\[
S(k,k') = \sum_{x=0}^{N-1} e^{i(\omega_k-\omega_{k'})x} = N\,\delta_{k,k'}
=\begin{cases} N, & k=k' \\ 0, & k\neq k', \end{cases}
\]
i.e.\ the discrete complex exponentials $e^{i\omega_k x}$ are orthogonal over one period.

\textbf{Step 4: Substitute back and conclude.}
Only the $k=k'$ terms of the double sum survive:
\[
\sum_{x=0}^{N-1} |f(x)|^2
= \frac{1}{N}\sum_{k=0}^{N-1}\sum_{k'=0}^{N-1} F(\omega_k)\,\overline{F(\omega_{k'})}\, N\,\delta_{k,k'}
= \sum_{k=0}^{N-1} F(\omega_k)\,\overline{F(\omega_k)}
= \sum_{k=0}^{N-1} |F(\omega_k)|^2,
\]
which is exactly Parseval's theorem, $\sum_{x} |f(x)|^2 = \sum_{\omega} |F(\omega)|^2$.
$\blacksquare$

\end{document}
